% !TEX root = missing_income.tex
% !TEX program = pdflatex
\documentclass[11pt]{article}

\usepackage[margin=1in]{geometry}
\usepackage{graphicx}
\usepackage{booktabs}
\usepackage{float}
\usepackage{hyperref}
\usepackage{caption}

\hypersetup{colorlinks=true, linkcolor=blue, citecolor=blue, urlcolor=blue}

\title{Characterizing Missing Income}
\author{Maria Jose Perez \and Juan Manuel Lozano \and Samuel Suárez Valle}
\date{\today}

\begin{document}

\maketitle

\noindent
The outcome variable, \texttt{y\_total\_m} (total monthly labor income), is missing for
1{,}778 of the 16{,}542 respondents in the analysis sample built in
\texttt{02\_Data\_Cleaning.r}. This document characterizes who these respondents are and
what that implies for handling the missingness in the rest of the analysis.

\section{Sample and a Filter Correction}

The analysis sample restricts the raw GEIH scrape to employed adults. An earlier version
of the filter used \texttt{dsi == 0} to mean ``employed,'' which looked reasonable but
was not equivalent: it let through 6{,}098 economically inactive respondents
(\texttt{ocu == 0}, \texttt{inac == 1}) --- respondents with no job to have missing income
\emph{from}. Under that filter, 77\% of the ``missing income'' group was simply the wrong
population, not a data problem.

Switching the filter to \texttt{age >= 18 \&\ ocu == 1} (occupied adults) drops the sample
from 22{,}640 to 16{,}542 rows and the missing-income rate from 34.8\% to
\textbf{10.75\%}. Every remaining row now has a full set of employment-module variables
(\texttt{relab}, \texttt{formal}, \texttt{sizeFirm}) --- the missingness left is a genuine
reporting pattern, not a sample-scope artifact.

\section{Demographic Balance}

\autoref{tab:balance} compares observations with and without \texttt{y\_total\_m} on age,
sex and education using two-sample t-tests. Age shows the clearest gap: respondents
missing income are on average 5 years older (44.0 vs.\ 38.9). The missing group is
somewhat less educated --- under-represented at ``secondary complete'' and over-represented
at ``none'' and ``primary incomplete'' --- and slightly more male (55.8\% vs.\ 52.7\%),
though both differences are small next to the pattern in \autoref{sec:relab}.

\begin{table}[!h]
\centering
\caption{\label{tab:balance}Balance table: missing vs. non-missing income}
\centering
\resizebox{\ifdim\width>\linewidth\linewidth\else\width\fi}{!}{
\begin{tabular}[t]{lrrrrr}
\toprule
Variable & Non-missing & Missing & Diff. & t-stat & p-value\\
\midrule
Age & 38.8882 & 43.9865 & -5.0983 & -13.7957 & 0.0000\\
Female & 0.4734 & 0.4421 & 0.0313 & 2.5095 & 0.0122\\
Male & 0.5266 & 0.5579 & -0.0313 & -2.5095 & 0.0122\\
None & 0.0066 & 0.0152 & -0.0086 & -2.8946 & 0.0038\\
Primary incomplete (1-4) & 0.0448 & 0.0602 & -0.0153 & -2.6027 & 0.0093\\
\addlinespace
Primary complete (5) & 0.0912 & 0.0922 & -0.0011 & -0.1466 & 0.8835\\
Secondary incomplete (6-10) & 0.1173 & 0.0917 & 0.0256 & 3.4938 & 0.0005\\
Secondary complete (11) & 0.3260 & 0.2655 & 0.0605 & 5.4181 & 0.0000\\
Tertiary & 0.4141 & 0.4753 & -0.0611 & -4.8805 & 0.0000\\
Private employee & 0.5931 & 0.3290 & 0.2641 & 22.2752 & 0.0000\\
\addlinespace
Government employee & 0.0387 & 0.0343 & 0.0044 & 0.9493 & 0.3426\\
Domestic worker & 0.0381 & 0.0084 & 0.0297 & 11.0735 & 0.0000\\
Self-employed & 0.2974 & 0.4021 & -0.1047 & -8.5664 & 0.0000\\
Employer & 0.0320 & 0.0861 & -0.0540 & -7.9331 & 0.0000\\
Unpaid family worker & 0.0000 & 0.1164 & -0.1164 & -15.3017 & 0.0000\\
\addlinespace
Unpaid worker (other household) & 0.0000 & 0.0231 & -0.0231 & -6.4764 & 0.0000\\
Day laborer & 0.0001 & 0.0000 & 0.0001 & 1.0000 & 0.3173\\
Other & 0.0005 & 0.0006 & 0.0000 & -0.0346 & 0.9724\\
Informal & 0.3976 & 0.5427 & -0.1452 & -11.6262 & 0.0000\\
Formal & 0.6024 & 0.4573 & 0.1452 & 11.6262 & 0.0000\\
\bottomrule
\end{tabular}}
\end{table}


\section{Employment Relationship}
\label{sec:relab}

The clearest driver of missingness is \texttt{relab}, the respondent's relationship to
their job. Standard salaried employment (``private employee'') makes up 59.3\% of
respondents with reported income but only 32.9\% of those missing it; self-employment and
business ownership run in the opposite direction (\autoref{tab:balance}).
\autoref{tab:relab-missing-rate} restates this as a rate rather than a composition: the
share of \emph{each} employment-relationship category that is itself missing income.

\begin{table}[H]
\centering
\caption{\label{tab:relab-missing-rate}Missing-income rate by employment relationship}
\centering
\resizebox{\ifdim\width>\linewidth\linewidth\else\width\fi}{!}{
\begin{tabular}[t]{lrrr}
\toprule
Employment relationship & N & N missing & Missing (\%)\\
\midrule
Private employee & 9342 & 585 & 6.3\\
Government employee & 632 & 61 & 9.7\\
Domestic worker & 578 & 15 & 2.6\\
Self-employed & 5106 & 715 & 14.0\\
Employer / business owner & 626 & 153 & 24.4\\
\addlinespace
Unpaid family worker & 207 & 207 & 100.0\\
Unpaid worker, another household's business & 41 & 41 & 100.0\\
Day laborer & 1 & 0 & 0.0\\
Other & 9 & 1 & 11.1\\
\bottomrule
\end{tabular}}
\end{table}


Unpaid family workers and unpaid workers in another household's business are missing
income for 100\% of respondents in those categories --- by construction, neither receives
a wage, so a missing \texttt{y\_total\_m} is the expected value for them, not
non-response. Self-employment (14.0\%) and, especially, employer/business-owner status
(24.4\%) show meaningfully elevated but far-from-total missingness, consistent with
income under-reporting rather than the absence of any income to report.

\section{Formality}

Consistent with the employment-relationship pattern, the missing-income group is more
informal: 45.7\% of respondents missing income are formal, versus 60.2\% of respondents
with reported income (\autoref{tab:balance}).

\section{Implications for Modeling}

Age, education and employment relationship are all observed variables. If an income
model already conditions on them, missingness can be treated as \textbf{missing at
random conditional on X}: dropping missing rows still identifies $E[Y \mid X]$, since
whatever predicts missingness is already in the model. This can be checked directly by
regressing an indicator for missingness on age, education and \texttt{relab} and
reporting how much of the variation it explains.

Where that defense is insufficient, two alternatives follow the same logic:

\begin{itemize}
  \item \textbf{Inverse-probability weighting} --- reweight observed rows by the inverse
  of their estimated probability of being observed (from the probit above), so the
  estimation sample looks like the full population again.
  \item \textbf{Multiple imputation} --- keep the full $n = 16{,}542$ rather than
  dropping 1{,}778 rows, at the cost of an explicit imputation-model assumption.
\end{itemize}

The two structurally-zero categories --- unpaid family and other-household workers --- are
a special case better handled on their own terms: their missing income reflects an
absent wage, not uncertainty to resolve.

\end{document}
